% !TEX program = xelatex
\documentclass[12pt,a4paper]{ctexart}  % 12pt → 正文小四号宋体; 题目默认小二黑体

% ============ 页面设置 ============
\usepackage[top=2.5cm, bottom=2.5cm, left=2.8cm, right=2.8cm]{geometry}
\usepackage{graphicx}
\usepackage{float}
\usepackage{hyperref}
\usepackage{amsmath,amssymb,bm}
\usepackage{booktabs}
\usepackage{caption}
\usepackage{enumitem}

\hypersetup{
    colorlinks=true,
    linkcolor=blue,
    citecolor=blue,
    urlcolor=blue
}

% ============ 标题格式（小二加黑）============
\ctexset{
    section/format = {\zihao{3}\heiti},
    subsection/format = {\zihao{4}\heiti},
}

% ============ 正文开始 ============
\begin{document}

% -------- 题目（小二加黑）--------
\begin{center}
    {\zihao{-2}\heiti 基于 DCGAN 的汉堡图像生成研究}

    \vspace{1em}
    {\zihao{4} 姓名：\underline{谢杭洋} \quad
     专业：\underline{食品科学与工程} \quad
     学号：\underline{24070001122}}
\end{center}

% ============ 摘要 ============
\vspace{1em}
{\zihao{4}\heiti 摘\ \ 要：}
针对小样本条件下结构化食品图像生成中存在的纹理模糊、边缘伪影和色彩失真等问题，
本文基于深度卷积生成对抗网络（DCGAN）架构，在963张手裁汉堡数据集上构建了128×128分辨率的图像生成模型。
模型采用转置卷积生成器与卷积判别器的经典结构，结合最小二乘生成对抗网络（LSGAN）损失、
单层特征匹配损失和Minibatch标准差正则化进行训练；
引入边界感知数据增强策略减轻白底填充引起的边缘伪影，并以Fr\'{e}chet Inception Distance（FID）作为定量评价指标。
实验结果表明，模型能够生成具有清晰多层堆叠结构的汉堡图像，生成器与判别器损失在训练过程中趋于动态均衡，FID值随训练推进持续下降。
引入判别器评分筛选机制后，进一步提升了输出图像的平均质量。
本文验证了经典DCGAN架构在小样本特定域食品图像生成中的有效性，为食品图像合成提供了一种可行方案。

\vspace{0.5em}
{\zihao{4}\heiti 关键词：}生成对抗网络；DCGAN；汉堡图像生成；FID指标；数据增强

\vspace{1em}
% ============ 1. 引言 ============
\section{引言}

生成对抗网络（Generative Adversarial Network, GAN）由Goodfellow等人\cite{goodfellow2014gan}于2014年首次提出，
其核心思想是通过生成器与判别器的对抗博弈，使生成器逐步学会从随机噪声中合成逼真样本。
经过十余年的发展，GAN已成为深度学习领域最具影响力的生成模型范式之一，
并已广泛应用于众多方向：
夏俊芳等\cite{xia2026cgan}将条件GAN应用于土壤养分高光谱数据增强，显著提升了小样本条件下的回归精度；
李云红等\cite{li2026catgan}提出融合坐标注意力机制的CAT-GAN，有效改善了文本到图像的跨模态生成质量；
张奥莹等\cite{zhang2026watermark}利用改进GAN实现图像水印检测，在降低参数量的同时达到99.58\%的检测准确率；
任帅等\cite{ren2026ganrf}将GAN优化随机森林模型用于地质灾害易发性评价，将AUC从0.897提升至0.953；
王强等\cite{wang2026survey}系统综述了数据稀缺场景下GAN高效图像生成的五大技术策略，
涵盖了域自适应、元学习、数据增强、架构优化和损失优化等方法。
上述应用表明，GAN技术已从最初的简单图像生成，拓展至数据增强、跨模态合成、安全检测、地学评估等多元化方向。

然而，在食品图像生成这一特定领域，相关研究工作仍相对有限，目前主要面临以下尚未解决的挑战：
\textbf{第一，结构保持困难。}食品图像（如汉堡）具有独特的多层堆叠结构（面包、肉饼、生菜、奶酪等），
生成模型在合成过程中容易丢失这种层次化的空间组织关系，导致生成图像出现结构扭曲或层次模糊。
\textbf{第二，纹理与边缘质量不足。}小样本条件下，传统DCGAN的转置卷积易产生棋盘状伪影和不均匀的重叠，
在图像边缘区域尤为明显，表现为白底填充边界处的高光和锯齿。
\textbf{第三，小数据集训练稳定性。}高质量的食品图像数据集通常规模较小（数百至数千张），
深度生成模型在样本量不足时容易出现模式崩溃、过拟合或训练不收敛等问题。

Radford等\cite{radford2016dcgan}提出的深度卷积生成对抗网络（DCGAN）建立了稳定的卷积生成架构基线，
通过批归一化和适当的激活函数设计，在多个自然图像数据集上取得了令人瞩目的生成效果。
本研究以DCGAN为核心架构，针对上述汉堡图像生成的具体挑战展开优化，
旨在探索小样本条件下结构化食品图像的生成方案，为食品设计可视化、菜单自动生成和合成数据增强等应用场景提供技术基础。

% ============ 2. 方法 ============
\section{方法}

\subsection{开发环境与工具}

本研究基于MATLAB R2024a平台开发\cite{matlab2024}，利用Deep Learning Toolbox
提供的\texttt{dlnetwork}、\texttt{dlarray}和\texttt{dlfeval}等自定义训练循环API构建模型；
使用Image Processing Toolbox进行图像预处理与后处理；
借助Parallel Computing Toolbox实现GPU加速训练（NVIDIA RTX 4060 8\,GB）；
采用Computer Vision Toolbox中的InceptionV3预训练网络提取特征用于FID计算。
代码框架参考了MathWorks官方提供的GAN花卉生成训练示例\cite{mathworksGAN}，
该示例展示了在MATLAB中实现自定义DCGAN训练循环的完整流程。

\subsection{生成器架构}

生成器$G$接收100维随机噪声向量$\mathbf{z} \sim \mathcal{N}(0, 1)$，通过6级转置卷积模块逐步上采样至128×128×3的RGB图像。
每级模块由转置卷积层（kernel size = 4）、批归一化和ReLU激活函数组成；
各级输出通道数依次为512、256、128、64、32、16，最后经卷积和tanh激活输出3通道彩色图像。
该设计遵循DCGAN的经典架构原则\cite{radford2016dcgan}：以转置卷积替代上采样加卷积的组合，以批归一化稳定各层激活分布，以ReLU加速收敛。

\subsection{判别器架构}

判别器$D$接收128×128×3图像，通过5级卷积下采样模块提取判别特征。
每级采用卷积层（kernel size = 4, stride = 2）、批归一化和LeakyReLU（$\alpha = 0.2$）激活，各级通道数为[64, 128, 256, 512, 512]。
网络末端引入Minibatch标准差层\cite{karras2018progressive}：计算批次内各空间位置的特征标准差，
将其作为一张额外的特征图拼接至特征向量，以增强判别器对生成样本多样性的感知能力，
从而抑制模式崩溃。最终通过全连接层输出标量判别分数。

\subsection{损失函数设计}

为缓解传统GAN交叉熵损失在训练后期出现的梯度消失问题，采用最小二乘生成对抗网络（LSGAN）损失\cite{mao2017lsgan}：

\begin{equation}
    \mathcal{L}_D = \frac{1}{2} \mathbb{E}_{\mathbf{x} \sim p_{\text{data}}}\left[(D(\mathbf{x}) - 1)^2\right]
    + \frac{1}{2} \mathbb{E}_{\mathbf{z} \sim p_{\mathbf{z}}}\left[D(G(\mathbf{z}))^2\right]
    \label{eq:lsgan_d}
\end{equation}

\begin{equation}
    \mathcal{L}_G = \frac{1}{2} \mathbb{E}_{\mathbf{z} \sim p_{\mathbf{z}}}\left[(D(G(\mathbf{z})) - 1)^2\right]
    \label{eq:lsgan_g}
\end{equation}

LSGAN以最小化Pearson $\chi^2$散度为目标，使生成样本向决策边界靠拢而非仅越过边界，
从而产生更高质量的生成图像\cite{mao2017lsgan}。

为进一步提升生成图像的纹理质量，引入特征匹配损失（Feature Matching）\cite{salimans2016improved}，
约束真实图像与生成图像在判别器中间层特征上的一致性（权重系数$\lambda_{\text{FM}} = 8$）：

\begin{equation}
    \mathcal{L}_{\text{FM}} = \left\| \mathbb{E}_{\mathbf{x} \sim p_{\text{data}}}[\phi(\mathbf{x})] - \mathbb{E}_{\mathbf{z} \sim p_{\mathbf{z}}}[\phi(G(\mathbf{z}))] \right\|_2^2
    \label{eq:fm}
\end{equation}

其中$\phi(\cdot)$表示判别器中间卷积层的特征映射。该损失从感知层面引导生成器匹配真实图像的统计特征分布。

\subsection{训练策略}

训练采用以下配置以保证收敛稳定性：（1）优化器使用Adam，判别器$\beta_1 = 0.0$、生成器$\beta_1 = 0.5$，
较低的判别器动量使其更快适应生成器的变化\cite{radford2016dcgan}；
（2）初始学习率$2\times10^{-4}$，采用余弦退火策略衰减至$2\times10^{-5}$，避免训练后期震荡；
（3）批量大小为32，在RTX 4060 8\,GB显存限制下平衡梯度估计质量与训练效率；
（4）生成器权重采用指数滑动平均（EMA，$\beta_{\text{EMA}} = 0.999$），推理时使用平滑后的权重以获得更稳定的输出；
（5）梯度裁剪阈值：判别器50、生成器25，防止单个异常样本造成梯度爆炸；
（6）边界感知数据增强：针对数据集中白底填充区域的边缘伪影问题，
在预处理阶段对图像边界区域注入随机噪声，迫使判别器学习忽略人工边框特征。

\subsection{定量评价指标}

采用Fr\'{e}chet Inception Distance（FID）\cite{heusel2017gans}作为生成质量的核心定量指标。
FID通过InceptionV3网络将真实图像与生成图像映射至特征空间，计算两者在该空间中的Fr\'{e}chet距离：

\begin{equation}
    \text{FID} = \left\| \bm{\mu}_r - \bm{\mu}_g \right\|_2^2
    + \operatorname{Tr}\left( \bm{\Sigma}_r + \bm{\Sigma}_g - 2(\bm{\Sigma}_r \bm{\Sigma}_g)^{1/2} \right)
    \label{eq:fid}
\end{equation}

其中$(\bm{\mu}_r, \bm{\Sigma}_r)$和$(\bm{\mu}_g, \bm{\Sigma}_g)$分别为真实图像与生成图像在InceptionV3全局平均池化层特征上的均值向量与协方差矩阵。
FID对模式崩溃和生成多样性不足敏感，值越低表示生成图像分布越接近真实图像分布\cite{heusel2017gans}。

此外，为从大量生成样本中自动筛选高质量输出，引入判别器评分筛选机制\cite{shmelkov2018good}：
训练完成后生成200张候选图像，逐张输入判别器进行评分，以sigmoid激活后的概率值作为质量分数，
保留分数最高的25张作为最终输出。该策略无需人工干预，有效提升了输出集合的平均质量。

% ============ 3. 实验 ============
\section{实验}

\subsection{数据来源与预处理}

本实验所用数据集为人工采集的汉堡正面照片集合（Burger\_v2），
共包含\textbf{963张}图像。原始图像来源于网络公开的美食图片和自摄照片，
涵盖了不同面包色泽、配料组合、拍摄角度和光照条件的汉堡样本，
具有一定的类内多样性。

数据预处理流程如下：
（1）手工裁剪去除背景杂乱元素，仅保留汉堡主体区域；
（2）对非正方形裁剪结果填充白色背景至统一的256×256像素，保持汉堡居于画幅中央；
（3）训练时统一缩放至128×128分辨率，以匹配生成器输入要求并降低计算开销；
（4）像素值归一化至$[-1, 1]$区间，与生成器输出层tanh激活函数的输出范围一致。

训练集使用全部963张图像，验证集采用训练过程中每5个epoch计算一次的FID指标。
每500次迭代自动保存一批生成样本至磁盘，用于训练过程的定性监控。

\begin{figure}[H]
    \centering
    \includegraphics[width=0.85\textwidth]{dataset_samples.png}
    \caption{Burger\_v2数据集样本展示（5×5拼贴图，共963张训练图像中的25张抽样）}
    \label{fig:dataset}
\end{figure}

图\ref{fig:dataset}展示了Burger\_v2数据集的25张抽样图像。
从图中可以看出，数据集具有以下特点：
（1）所有图像均经手工裁剪并填充白底至256×256像素，汉堡主体居中且占比一致；
（2）面包色泽覆盖浅金至深褐多种烘焙程度，配料组合（肉饼、生菜、奶酪、酱料）多样；
（3）拍摄角度以正上方俯拍为主，光照条件涵盖自然光和室内灯光。
该数据集规模为963张，属于典型的小样本特定域图像数据集，对生成模型的训练稳定性和泛化能力提出了挑战。

\subsection{生成图像结果}

\begin{figure}[H]
    \centering
    \includegraphics[width=0.85\textwidth]{montage_5x5.png}
    \caption{DCGAN生成的5×5汉堡图像拼贴图（经D判别器评分筛选后的Top-25输出，输出分辨率为512×512）}
    \label{fig:montage}
\end{figure}

图\ref{fig:montage}展示了模型经D评分筛选后输出的25张最佳汉堡图像的5×5拼贴图。
从生成结果可以看出：
\textbf{在结构方面，}汉堡标志性的水平多层堆叠特征得到了较好的保留——上层圆形面包、中层肉饼与蔬菜、下层面包的层次关系清晰可辨，
未出现层次融合或结构崩溃现象，验证了DCGAN架构对结构化物体的建模能力。
\textbf{在多样性方面，}不同生成样本在面包色泽（浅金色至深褐色）、配料排列（生菜的位置与厚度）和整体构图上呈现明显差异，
说明模型未陷入单一模式的重复生成。
\textbf{在色彩方面，}生成图像的整体色调在汉堡食材的自然色彩范围内（面包的棕黄色系、蔬菜的绿色点缀、肉饼的深褐色），
未出现严重的色彩溢出或色块异常。

\subsection{训练过程分析}

\begin{figure}[H]
    \centering
    \includegraphics[width=0.92\textwidth]{result.png}
    \caption{训练过程曲线图（包含生成器/判别器损失曲线、FID变化趋势、生成器得分分布等训练监控指标）}
    \label{fig:training}
\end{figure}

图\ref{fig:training}展示了完整训练过程中的各项监控指标变化曲线。训练曲线揭示了以下规律：

（1）\textbf{损失收敛性：}生成器与判别器的LSGAN损失在训练初期经历了典型的对抗博弈振荡阶段，
约在第50个epoch后逐步趋于动态均衡——判别器损失维持在0.1--0.3的稳定区间，
生成器损失则在0.5--1.0范围内小幅波动，未出现损失发散或单方压倒性优势，
验证了LSGAN损失函数在两方博弈中的稳定特性\cite{mao2017lsgan}。

（2）\textbf{FID下降趋势：}FID指标随训练推进整体呈现持续下降的态势，
从训练初期的较高值（约300--400）逐步改善至更低的稳定水平，
表明生成图像的分布与真实图像分布在InceptionV3特征空间中逐步接近。
FID的下降与生成图像视觉质量的改善呈正相关，验证了该指标在训练过程中作为自动评价标准的有效性。

（3）\textbf{训练鲁棒性：}整个训练过程中各项指标未出现NaN或Inf等异常数值，
未发生模式崩溃（表现为生成图像突然丧失多样性），
验证了梯度裁剪、EMA和余弦退火学习率等训练策略在维持训练稳定性方面的协同作用。

% ============ 4. 讨论 ============
\section{讨论}

\subsection{方法的优势}

本方法在小样本食品图像生成任务中展现出了以下优势：
\textbf{架构简洁有效：}经典DCGAN架构在仅963张训练图像的条件下，
成功学习到了汉堡的多层结构特征，证明在小样本特定域任务中，
经过适当调优的经典架构其表现不亚于复杂模型。
\textbf{训练稳定性：}LSGAN损失配合特征匹配损失和Minibatch标准差正则化，
在整个训练过程中保持了良好的对抗平衡，未发生模式崩溃——这对于小样本GAN训练尤为重要。
\textbf{筛选机制实用：}利用训练好的判别器进行自动评分筛选，
在无人工标注的情况下有效提升了最终输出图像的平均质量，
为实际应用中的批量生成与自动择优提供了一种低成本方案。
\textbf{工程可复现：}全部基于MATLAB官方工具箱实现，代码框架参考官方示例，
具有良好的可复现性和教学参考价值。

\subsection{不足与局限}

同时，本研究存在以下不足：
\textbf{纹理锐度不足：}部分生成图像在局部细节区域（如芝麻颗粒、生菜叶片边缘）呈现柔和的失焦感，
这与转置卷积的固有棋盘效应以及小数据集对高频细节学习的不足有关。
Salimans等\cite{salimans2016improved}也指出，单层特征匹配虽然能提升全局结构，
但对局部纹理的约束力有限。
\textbf{光照一致性欠佳：}不同生成样本之间的光源方向和高光分布存在不一致，
部分原因是训练数据本身由不同光照条件下的照片组成，而模型缺乏显式的光照建模机制。
\textbf{分辨率的限制：}受限于RTX 4060的8\,GB显存，模型在128×128分辨率下进行训练，
直接生成的图像在放大至展示尺寸时细节有所不足，需要依赖超分辨率后处理进行弥补。

\subsection{可能的改进方向}

针对上述不足，后续研究可从以下方向进行改进：
（1）引入ESRGAN等超分辨率模型\cite{wang2018esrgan}对生成的128×128图像进行后处理增强至512×512，
在不增加GAN训练开销的前提下提升输出分辨率和纹理细节；
（2）采用StyleGAN风格的映射网络和自适应实例归一化（AdaIN），
通过分离式潜在空间控制增强生成图像的多样性和可控性；
（3）引入R1梯度惩罚\cite{karras2018progressive}替代梯度裁剪，
进一步提高判别器的训练平滑度和生成图像质量；
（4）扩大数据集规模，结合更丰富的在线数据增强策略（随机翻转、色彩抖动、随机裁剪），
提升模型的泛化能力；
（5）探索扩散模型（Diffusion Models）在该数据集上的生成效果，与GAN进行定性和定量的对比分析。

% ============ 5. 结论 ============
\section{结论}

本文针对小样本条件下结构化食品图像生成任务，
以DCGAN为基线架构，在963张手裁汉堡图像数据集上构建了128×128分辨率的图像生成模型。
通过LSGAN损失函数解决传统GAN的梯度消失问题，
以单层特征匹配损失从感知层面约束生成图像的纹理质量，
以Minibatch标准差层增强判别器的多样性感知以抑制模式崩溃，
并结合边界感知数据增强和余弦退火学习率等训练策略保持训练的稳定收敛。

实验结果表明：模型成功学习到了汉堡标志性的水平多层堆叠结构，
生成图像在结构完整性、色彩合理性和样本多样性方面均达到了预期水平；
FID指标随训练推进持续下降，定量验证了生成质量的有效改善；
训练全程未出现模式崩溃或数值发散，验证了本文训练策略组合的鲁棒性。
判别器自动评分筛选机制的引入，进一步在不依赖人工判断的条件下提升了最终输出质量。

本研究验证了经典DCGAN架构在小样本特定域食品图像生成任务中的可行性与有效性。
在回顾整个研究过程的基础上，提出以下后续研究方向：
（1）将超分辨率模型与GAN级联，在不增加GAN训练开销的前提下提升输出分辨率和细节；
（2）引入StyleGAN风格的风格调制机制，提升生成结果的多样性和可控性；
（3）探索扩散模型在该任务上的潜力，与GAN方案进行系统性对比；
（4）扩充数据集并引入更丰富的数据增强策略，进一步提升模型泛化能力。

% ============ 参考文献 ============
\begin{thebibliography}{99}
\raggedright

\bibitem{goodfellow2014gan}
Goodfellow I, Pouget-Abadie J, Mirza M, et al.
Generative adversarial nets[C].
Advances in Neural Information Processing Systems, 2014, 27: 2672--2680.

\bibitem{radford2016dcgan}
Radford A, Metz L, Chintala S.
Unsupervised representation learning with deep convolutional generative adversarial networks[C].
International Conference on Learning Representations, 2016.

\bibitem{mao2017lsgan}
Mao X, Li Q, Xie H, et al.
Least squares generative adversarial networks[C].
Proceedings of the IEEE International Conference on Computer Vision, 2017: 2794--2802.

\bibitem{salimans2016improved}
Salimans T, Goodfellow I, Zaremba W, et al.
Improved techniques for training GANs[C].
Advances in Neural Information Processing Systems, 2016, 29: 2234--2242.

\bibitem{karras2018progressive}
Karras T, Aila T, Laine S, et al.
Progressive growing of GANs for improved quality, stability, and variation[C].
International Conference on Learning Representations, 2018.

\bibitem{heusel2017gans}
Heusel M, Ramsauer H, Unterthiner T, et al.
GANs trained by a two time-scale update rule converge to a local Nash equilibrium[C].
Advances in Neural Information Processing Systems, 2017, 30: 6626--6637.

\bibitem{shmelkov2018good}
Shmelkov K, Schmid C, Alahari K.
How good is my GAN?[C].
Proceedings of the European Conference on Computer Vision, 2018: 213--229.

\bibitem{wang2018esrgan}
Wang X, Yu K, Wu S, et al.
ESRGAN: Enhanced super-resolution generative adversarial networks[C].
Proceedings of the European Conference on Computer Vision Workshops, 2018.

\bibitem{xia2026cgan}
夏俊芳, 苑宏文, 张浩, 等.
基于条件生成对抗网络数据增强的土壤总氮高光谱反演[J].
农业工程学报, 2026, 42(5): 137--146.

\bibitem{li2026catgan}
李云红, 张琪琪, 陈锦妮, 等.
基于生成对抗网络和坐标注意力机制的文本生成图像算法[J].
浙江大学学报(工学版), 2026, 60(6): 1213--1220.

\bibitem{zhang2026watermark}
张奥莹, 王树梅, 顾张清.
基于改进生成对抗网络的图像水印检测算法[J].
微电子学与计算机, 2026, 43(4): 50--65.

\bibitem{ren2026ganrf}
任帅, 李有淋, 余杰.
基于生成对抗网络优化随机森林模型的西藏定结县地质灾害易发性评价[J].
中国地质灾害与防治学报, 2026, 37(2): 156--168.

\bibitem{wang2026survey}
王强, 蒋国权, 翁年凤, 等.
基于生成对抗网络的高效图像生成技术综述[J/OL].
计算机科学, 2026, 1--29. \url{https://link.cnki.net/urlid/50.1075.tp.20260526.1655.002}.

\bibitem{matlab2024}
The MathWorks Inc.
MATLAB (R2024a)[CP/OL].
Natick, Massachusetts, 2024. \url{https://www.mathworks.com/products/matlab.html}.

\bibitem{mathworksGAN}
The MathWorks Inc.
Train generative adversarial network (GAN)[EB/OL].
Deep Learning Toolbox Documentation, 2024. \url{https://www.mathworks.com/help/deeplearning/ug/train-generative-adversarial-network.html}.

\end{thebibliography}

\end{document}
